\documentclass[12pt,a4paper,UTF8]{ctexart}
\usepackage[a4paper,top=22mm,bottom=22mm,left=22mm,right=22mm,headheight=15pt]{geometry}
\usepackage{amsmath,amssymb,mathtools,bm}
\usepackage{graphicx}
\usepackage{xcolor}
\usepackage{booktabs,longtable,array}
\usepackage{enumitem}
\usepackage{fancyhdr}
\usepackage{caption}
\usepackage{float}
\usepackage{microtype}
\usepackage{setspace}
\usepackage[hidelinks,unicode]{hyperref}
\usepackage[most]{tcolorbox}

\definecolor{ContestBlue}{HTML}{17365D}
\definecolor{SoftBlue}{HTML}{EEF4FA}
\definecolor{SoftGray}{HTML}{F5F6F7}
\hypersetup{pdftitle={第97套},pdfauthor={物理竞赛教学资料}}

\setstretch{1.24}
\setlength{\parindent}{2em}
\setlength{\parskip}{0.35em}
\setlist{itemsep=0.35em,topsep=0.4em}
\captionsetup{font=small,labelfont=bf,justification=centering}
\setcounter{secnumdepth}{0}

\pagestyle{fancy}
\fancyhf{}
\fancyhead[L]{\small 第97套}
\fancyfoot[C]{\thepage}
\renewcommand{\headrulewidth}{0.4pt}

\ctexset{
  section={format=\Large\bfseries\color{ContestBlue},beforeskip=1.2em,afterskip=0.7em},
  subsection={format=\large\bfseries\color{ContestBlue},beforeskip=1.0em,afterskip=0.5em},
  subsubsection={format=\normalsize\bfseries,beforeskip=0.8em,afterskip=0.35em}
}

\makeatletter
\newcommand{\tightlist}{\setlength{\itemsep}{0pt}\setlength{\parskip}{0pt}}
\newcommand{\maxwidth}{\ifdim\Gin@nat@width>\linewidth\linewidth\else\Gin@nat@width\fi}
\newcommand{\maxheight}{\ifdim\Gin@nat@height>0.72\textheight 0.72\textheight\else\Gin@nat@height\fi}
\newcommand{\pandocbounded}[1]{\begingroup\setkeys{Gin}{width=\maxwidth,height=\maxheight,keepaspectratio}#1\endgroup}
\makeatother

\begin{document}

\begin{titlepage}
  \centering
  \vspace*{58mm}
  {\Huge\bfseries\color{ContestBlue} 第97套\par}
  \vfill
  \vspace*{28mm}
\end{titlepage}


\section{第 1 题}\label{ux7b2c-1-ux9898}

设固定光滑椭圆的方程为

\[\frac{x^2}{a^2}+\frac{y^2}{b^2}=1.\]

\begin{enumerate}
\def\labelenumi{\arabic{enumi}.}
\tightlist
\item
  一足够长的直角刚性轻架搭在椭圆上，并可在 \(x\)--\(y\)
  平面内无摩擦滑动。轻架的直角顶点 \(M\) 处固定一质量为 \(m\)
  的质点。初始时 \(M\) 静止在其轨迹的最高点，受到微小扰动后开始运动。求
  \(OM\) 方向转过多大角度时，轻架开始与椭圆分离。
\item
  将同一椭圆置于由光滑水平地面和光滑竖直墙组成的直角墙角内，使椭圆同时与地面、墙面相切；初始时椭圆中心位于
  \((b,a)\)，且椭圆的全部质量集中在中心。椭圆受扰动后逆时针转动，并在保持两处相切的条件下滑动。若要求椭圆中心到达
  \((a,b)\) 以前两处接触均不失效，求 \(a,b\) 应满足的条件。
\end{enumerate}

\begin{figure}[H]
\centering
\includegraphics[width=0.9\linewidth,height=\textheight,keepaspectratio,alt={图 1-1　几何约束：直角切线架的顶点以及同时与两坐标轴相切的椭圆。}]{../figures/question-01.pdf}
\caption*{图
1-1　几何约束：直角切线架的顶点以及同时与两坐标轴相切的椭圆。}
\end{figure}

\newpage

\section{第 2 题}\label{ux7b2c-2-ux9898}

质量为 \(m\) 的小球以速度 \(v_0\) 斜向下撞击一块初始静止、质量同为 \(m\)
的水平木板左端。入射速度与竖直方向的夹角为
\(\theta\)，小球与木板间的动摩擦因数为
\(\mu\)，木板与地面间光滑。每次碰撞中，小球的竖直速度分量均等大反向。重力加速度为
\(g\)。

\begin{enumerate}
\def\labelenumi{\arabic{enumi}.}
\tightlist
\item
  暂将小球视为质点，求第一次碰撞后瞬间小球和木板的速度。
\item
  取 \(\theta=45^\circ\)、\(\mu=0.1\)、\(v_0=10\,\mathrm{m/s}\)、木板长
  \(l=7\,\mathrm m\)、\(g=10\,\mathrm{m/s^2}\)。通过计算判断小球最终是否会从木板上掉下；若会，求使其不掉下所需的最小板长。
\item
  将小球视为半径为 \(r\) 的均匀实心球，转动惯量
  \(I=2mr^2/5\)。仍取第（2）问数据，判断小球落到木板所在高度以前，是否会与木板发生第二次碰撞。
\end{enumerate}

\begin{figure}[H]
\centering
\includegraphics[width=0.9\linewidth,height=\textheight,keepaspectratio,alt={图 2-1　小球斜碰可水平移动的木板。}]{../figures/question-02.pdf}
\caption*{图 2-1　小球斜碰可水平移动的木板。}
\end{figure}

\newpage

\section{第 3 题}\label{ux7b2c-3-ux9898}

设整个空间充满均匀磁化的永磁介质，磁化强度为

\[\bm M=\frac{\sum \bm m}{V},\]

其中 \(\bm m\) 为单个磁矩，\(V\) 为体积。在介质中挖出一个空腔，且
\(\bm M\) 与空腔的对称轴平行。求空腔中心的磁感应强度 \(\bm B\)
与磁场强度 \(\bm H\)。分别考虑：

\begin{enumerate}
\def\labelenumi{\arabic{enumi}.}
\tightlist
\item
  厚度为 \(h\)、半径为 \(R\) 且 \(R\gg h\) 的薄圆盘形空腔；
\item
  半径为 \(r\)、长度为 \(l\) 且 \(l\gg r\) 的细长圆柱形空腔；
\item
  半径为 \(R\) 的球形空腔。
\end{enumerate}

\begin{figure}[H]
\centering
\includegraphics[width=0.9\linewidth,height=\textheight,keepaspectratio,alt={图 3-1　轴线与磁化强度平行的薄圆盘形、细长圆柱形和球形空腔。}]{../figures/question-03.pdf}
\caption*{图 3-1　轴线与磁化强度平行的薄圆盘形、细长圆柱形和球形空腔。}
\end{figure}

\newpage

\section{第 4 题}\label{ux7b2c-4-ux9898}

在光滑水平桌面内建立 \(x\)--\(y\) 坐标系，并施加竖直向下、大小为 \(B\)
的匀强磁场。一质量为 \(m\)、电荷量为 \(+q\) 的点电荷从原点 \(O\) 沿
\(+x\) 方向以初速度 \(v_0\) 运动。

\begin{enumerate}
\def\labelenumi{\arabic{enumi}.}
\tightlist
\item
  不加其他装置时，求点电荷的轨迹。
\item
  在 \(O\) 点与点电荷之间连接一根劲度系数为
  \(k\)、原长可忽略不计的绝缘弹簧。仍使点电荷从 \(O\) 点沿 \(+x\) 方向以
  \(v_0\) 运动，求轨迹方程。
\end{enumerate}

\begin{figure}[H]
\centering
\includegraphics[width=0.706\linewidth,height=\textheight,keepaspectratio,alt={图 4-1　竖直向下磁场中，以原点为固定端的零原长弹簧连接正电荷。}]{../figures/question-04.pdf}
\caption*{图 4-1　竖直向下磁场中，以原点为固定端的零原长弹簧连接正电荷。}
\end{figure}

\newpage

\section{第 5 题}\label{ux7b2c-5-ux9898}

一竖直光滑圆锥的半顶角为 \(\theta\)。圆锥内固定一半径为 \(r_0\)
的水平圆形导线环，其下方放置一质量为 \(m\)
的均匀导电弹性环。弹性环原长可忽略不计，平衡时保持水平，半径为
\(R_0\)，且 \(R_0\gg r_0\)；重力加速度为 \(g\)。初始平衡时大环中无电流。

\begin{enumerate}
\def\labelenumi{\arabic{enumi}.}
\tightlist
\item
  求弹性环的劲度系数 \(k\)；
\item
  求平衡位置两环的互感系数 \(M\)；
\item
  若小环中通有恒定电流 \(I_0\)，求穿过大环的外磁通 \(\Phi_0\)；
\item
  大环电阻为 \(R_a\)、自感可忽略。给大环一个沿圆锥母线方向、大小为
  \(v_0\) 的小速度，求此时大环质心的加速度；
\item
  大环自感为
  \(L\)、电阻可忽略。给大环一个沿母线方向的小扰动，并设其始终保持水平，求小振动周期
  \(T\)。
\end{enumerate}

\begin{figure}[H]
\centering
\includegraphics[width=0.752\linewidth,height=\textheight,keepaspectratio,alt={图 5-1　竖直圆锥内的固定载流小环与可沿母线运动的导电弹性大环。}]{../figures/question-05.pdf}
\caption*{图 5-1　竖直圆锥内的固定载流小环与可沿母线运动的导电弹性大环。}
\end{figure}

\newpage

\section{第 6 题}\label{ux7b2c-6-ux9898}

一理想容器总体积为 \(v_0\)，由隔板分成两个等体积区域。左室有 \(N\)
个称为``人''的质点，右室有 \(N\) 个称为``妖''的质点；\(N\)
足够大，两边均可视为单原子理想气体。

\begin{enumerate}
\def\labelenumi{\arabic{enumi}.}
\tightlist
\item
  两边初温均为
  \(T_0\)，且``人''\,``妖''可分辨。抽去隔板并达到均匀混合，求熵变。
\item
  左、右初温分别为
  \(T_0,T_1\)，且两类粒子可分辨。抽去隔板并达到平衡，求熵变。
\item
  左、右初温仍为
  \(T_0,T_1\)，但两边粒子完全不可分辨。抽去隔板并达到平衡，求熵变。
\item
  初态左、右两边分别由``人--人''和``妖--妖''双原子分子组成，温度均为
  \(T_0\)；``人''和``妖''除标签外物理、化学性质相同，但视为两类粒子。沿用题首记号，每侧仍有
  \(N\) 个原子，因而每侧初始有 \(N/2\)
  个双原子分子。混合后还可形成``人--妖''分子。求反应平衡后三类分子的比例及全过程熵变。
\end{enumerate}

\begin{figure}[H]
\centering
\includegraphics[width=0.869\linewidth,height=\textheight,keepaspectratio,alt={图 6-1　等体积两室内可分辨或不可分辨粒子的混合。}]{../figures/question-06.pdf}
\caption*{图 6-1　等体积两室内可分辨或不可分辨粒子的混合。}
\end{figure}

\newpage

\section{第 7 题}\label{ux7b2c-7-ux9898}

一厚度为 \(l\) 的单轴晶体波片，其表面法线与光轴均沿 \(y\)
轴。一束真空波长为 \(\lambda\) 的平行光在 \(x\)--\(y\)
平面内入射，入射方向与 \(y\) 轴夹角为
\(\theta\)。入射光为线偏振光，垂直于入射面的 \(s\) 分量与平行于入射面的
\(p\) 分量等振幅、同相。普通光折射率为 \(n_o\)；当非常光波矢与光轴夹角为
\(\alpha\) 时，折射率满足

\[\frac1{n^2(\alpha)}=\frac{\cos^2\alpha}{n_o^2}+\frac{\sin^2\alpha}{n_e^2}.\]

\begin{enumerate}
\def\labelenumi{\arabic{enumi}.}
\tightlist
\item
  判断 \(s\) 光对应 \(o\) 光还是 \(e\) 光，并求 \(o\) 光折射角 \(i_o\)；
\item
  根据惠更斯原理画出 \(e\) 光从原点进入晶体后，经 \(\Delta t\)
  形成的椭圆子波面；
\item
  由空气中波前与椭圆子波面的公切线求 \(e\) 光波矢折射角 \(i_e\)；
\item
  两束光出射后形成平行光，求其光程差 \(\Delta\)。
\end{enumerate}

\begin{figure}[H]
\centering
\includegraphics[width=0.547\linewidth,height=\textheight,keepaspectratio,alt={图 7-1　光轴垂直表面的单轴晶体中，普通光与非常光的折射。}]{../figures/question-07.pdf}
\caption*{图 7-1　光轴垂直表面的单轴晶体中，普通光与非常光的折射。}
\end{figure}

\newpage

\section{第 8 题}\label{ux7b2c-8-ux9898}

在 \(x\)--\(y\) 平面内有一半径为 \(r\) 的半圆，其直径与 \(y\)
轴重合。半圆弧上均匀分布 \(N\)
个相同、彼此独立且初始静止的小球，总质量为 \(M\)，相邻小球对应的圆心角为
\(\pi/(N-1)\)。另一质量为 \(m\) 的小球以速度 \(v_0\) 入射，其方向与
\(x\) 轴成 \(\pi/[2(N-1)]\) 的向下夹角。在适当条件下，入射球依次与 \(N\)
个静止小球发生完全弹性碰撞，最后以速度 \(v_1\) 沿关于 \(x\)
轴对称的方向反弹。

\begin{enumerate}
\def\labelenumi{\arabic{enumi}.}
\tightlist
\item
  在 \(v_0\ll c\) 的非相对论条件下，求使上述过程能够发生的 \(M/m\)
  最小值；
\item
  取到最小质量比时，求 \(v_1/v_0\)； 3.1 当 \(v_0\)
  不可忽略相对论效应，且题中质量均为静止质量时，重新求第（1）问； 3.2
  相对论条件下第（2）问是否有简单显式表达式？若没有，给出最终速度的计算方法。
\end{enumerate}

\begin{figure}[H]
\centering
\includegraphics[width=0.564\linewidth,height=\textheight,keepaspectratio,alt={图 8-1　入射球沿半圆球列连续偏转并关于 x 轴对称反弹。}]{../figures/question-08.pdf}
\caption*{图 8-1　入射球沿半圆球列连续偏转并关于 \(x\) 轴对称反弹。}
\end{figure}

\end{document}
